% Work-MAT | 2. rocnik | pripravovana hodina 24
% Návaznost: uživatelské připravy.zip (hodiny 01–20), autorská učebnice na GitHubu.
% Realisticky.cz: didaktická inspirace v posloupnosti a typech úloh; zadání jsou vlastní.
\documentclass[11pt,a4paper]{article}
\input{../../sablony/pripravy/priprava.sty}
\nastavPripravu{2. ročník}{Exponenciální rovnice}{24}{Smíšené substituce}
\begin{document}
\blok{Motivace a návaznost}
Dnešní úlohy začínají rozhodováním: nejprve převeď výrazy na jeden základ, až potom zvol substituci. Někdy po převodu vyjde kvadratická rovnice.
\poznamka{Orientační nabídka pro 45--55 minut: krátké zopakování a motivace, nový výklad, tři společné vzorové úlohy, samostatná práce. Vyber podle tempa třídy; zbývající zadání slouží jako rezerva.}
\blok{Výklad}
Volíme $t=a^x>0$ a každý člen rovnici přepíšeme pomocí $t$. Po vyřešení pomocné rovnice vždy provedeme zpětnou substituci.
\blok{Společné příklady}
\begin{prikladReseniBox}[title={Příklad 1}]
Vyřeš, resp. urči: $16^x-5\cdot4^x+4=0$.
\tcblower
\[\begin{aligned}t=4^x>0,\quad t^2-5t+4=0 \\ t\in\{1,4\} \\ 4^x=1\;\lor\;4^x=4\;\Longrightarrow\;K=\{0,1\}\end{aligned}\]
\end{prikladReseniBox}
\begin{prikladReseniBox}[title={Příklad 2}]
Vyřeš, resp. urči: $9^x-4\cdot3^x+3=0$.
\tcblower
\[\begin{aligned}t=3^x>0,\quad (t-1)(t-3)=0 \\ K=\{0,1\}\end{aligned}\]
\end{prikladReseniBox}
\begin{prikladReseniBox}[title={Příklad 3}]
Vyřeš, resp. urči: $2^{2x+2}-5\cdot2^{x+1}+4=0$.
\tcblower
\[\begin{aligned}t=2^{x+1}>0,\quad (t-1)(t-4)=0 \\ x+1=0\;\lor\;x+1=2 \\ K=\{-1,1\}\end{aligned}\]
\end{prikladReseniBox}
\poznamka{Na tabuli porovnej možné substituce $t=2^x$ a $t=2^{x+1}$ a ukaž, proč někdy druhá zkrátí zápis.}
\blok{Samostatné procvičování}
\begin{zadaniReseniBox}[title={Samostatná práce 1}]
\begin{enumerate}[label=\alph*)]
\item $4^x-5\cdot2^x+4=0$.
\item $16^x-5\cdot4^x+4=0$.
\item $9^x-4\cdot3^x+3=0$.
\item $25^x-6\cdot5^x+5=0$.
\item $4^{x+1}-5\cdot2^{x+1}+4=0$.
\item Je $t=a^x=-4$ přípustné?.
\end{enumerate}
\tcblower
\begin{enumerate}[label=\alph*)]
\item $K=\{0,2\}$.
\item $K=\{0,1\}$.
\item $K=\{0,1\}$.
\item $K=\{0,1\}$.
\item $K=\{-1,1\}$.
\item Není..
\end{enumerate}
\end{zadaniReseniBox}
\blok{Dobrovolné rozšíření}
\blok{Úkol na doma}
Při početních úlohách napiš postup, zvolenou metodu a zkoušku v původním zadání. U soustav kontroluj obě rovnice; u nerovnic výsledek zakresli na osu.
\begin{zadaniReseniBox}[title={Domácí úkol}]
\begin{enumerate}[label=\alph*)]
\item $9^x+2\cdot3^x-15=0$.
\item $16^x-5\cdot4^x+4=0$.
\item $4^x-6\cdot2^x+8=0$.
\item $25^x-6\cdot5^x+5=0$.
\item $2^{2x+2}-5\cdot2^{x+1}+4=0$.
\item $9^x-10\cdot3^x+9=0$.
\item Zdůvodni podmínku $t>0$..
\end{enumerate}
\tcblower
\begin{enumerate}[label=\alph*)]
\item $K=\{1\}$.
\item $K=\{0,1\}$.
\item $K=\{1,2\}$.
\item $K=\{0,1\}$.
\item $K=\{-1,1\}$.
\item $K=\{0,2\}$.
\item $Každá exponenciální hodnota je kladná.$.
\end{enumerate}
\end{zadaniReseniBox}
\poznamka{Domácí úkol obsahuje 7 položek; předpokládaná domácí práce přibližně 15--30 minut. Po větších blocích jsou vědomě promíchané dříve probrané typy.}
\end{document}
